
\chapter*{Glossaire}
\addcontentsline{toc}{chapter}{Glossaire}

\begin{description}[leftmargin=3.2cm,style=nextline]
  \item[Agent] Utilisateur chargé d'importer le fichier annuel des consommations dans l'application.
  
  \item[Contrat] Association entre un client, un quartier et un type d'activité, définissant les conditions de gestion de ses consommations.
  
  \item[CSRF] (Cross-Site Request Forgery) Falsification de requête intersite. Cette attaque consiste à faire exécuter à un utilisateur authentifié une action qu'il n'a pas volontairement demandée. La protection repose notamment sur l'utilisation de jetons CSRF.
  
  \item[Facture] Document généré à partir d'une consommation validée, en appliquant les tarifs correspondants et la TVA.
  
  \item[IHM] Interface homme-machine. Ensemble des éléments permettant à l'utilisateur d'interagir avec l'application.
  
  \item[MVC] (Modèle--Vue--Contrôleur) Architecture logicielle qui sépare la gestion des données, leur présentation et le traitement des actions de l'utilisateur.
  
  \item[PDO] (PHP Data Objects) Extension PHP permettant d'accéder à différentes bases de données à travers une interface commune.
  
  \item[Quartier] Zone géographique associée à un fournisseur et à un type d'activité, utilisée pour organiser les clients et leurs contrats.
  
  \item[Réclamation] Demande formulée par un client concernant un problème lié à son service, notamment une fuite, une facture ou une autre situation.
  
  \item[TVA] Taxe sur la valeur ajoutée. Dans le cadre du projet, son taux est fixé à 18~\%.
  
  \item[XSS] (Cross-Site Scripting) Vulnérabilité permettant d'injecter des scripts malveillants dans une page web consultée par d'autres utilisateurs.
\end{description}